% =================================================
% Справочный лист
% =================================================

\begin{center}
  {\fontsize{25}{30}\selectfont\bfseries\sffamily\color{Primary}
    Ломаная и многоугольник}

  \vspace{0.3em}

  {\large\sffamily\color{GrayText}
    Звенья, вершины, стороны и диагонали}
\end{center}

\vspace{0.45em}

\section{Ломаная}

\openbrokenfigure

\begin{propertybox}[Определение]
  \textbf{Ломаная} состоит из отрезков, последовательно соединённых своими
  концами. Эти отрезки называют \textbf{звеньями}, а их концы ---
  \textbf{вершинами ломаной}.
\end{propertybox}

Ломаную называют, перечисляя её вершины по порядку. На рисунке изображена
ломаная \(ABCDE\). Она содержит четыре звена:
\[
  AB,\quad BC,\quad CD,\quad DE.
\]

\begin{warningbox}[Не пропускайте вершины]
  Название должно передавать путь вдоль ломаной. Запись \(ABDE\) не подходит:
  она пропускает вершину \(C\).
\end{warningbox}

\section{Замкнутая и незамкнутая ломаные}

\openclosedfigure

У незамкнутой ломаной два разных конца. У замкнутой ломаной последний конец
соединён с первым.

\begin{notebox}[Самопересечение]
  Замкнутая ломаная может пересекать сама себя. Для границы многоугольника
  используют замкнутую ломаную \textbf{без самопересечений}.
\end{notebox}

\section{Многоугольник}

\polygonpartsfigure

\begin{propertybox}[Определение]
  Часть плоскости, ограниченная замкнутой ломаной без самопересечений,
  называется \textbf{многоугольником}. Звенья границы называют
  \textbf{сторонами}, а их концы --- \textbf{вершинами многоугольника}.
\end{propertybox}

У многоугольника число сторон равно числу вершин.

\polygonnamesfigure

\begin{center}
  \renewcommand{\arraystretch}{1.35}
  \begin{tabular}{c c c c}
    \toprule
    Число сторон & 3 & 4 & 5 \\
    \midrule
    Название & треугольник & четырёхугольник & пятиугольник \\
    \bottomrule
  \end{tabular}
\end{center}

\section{Как правильно назвать многоугольник}

\correctnamingfigure

Вершины перечисляют \textbf{подряд по границе} в любом направлении. Начать
можно с любой вершины. Например, один и тот же многоугольник можно назвать
\(ABCDEF\), \(BCDEFA\) или \(AFEDCB\).

\begin{warningbox}[Название не является произвольным набором букв]
  В записи \(ACBDEF\) после \(A\) указана несоседняя вершина \(C\), поэтому
  такая запись не называет изображённый многоугольник.
\end{warningbox}

\clearpage
\section{Соседние вершины и диагонали}

Две вершины называют \textbf{соседними}, если они являются концами одной
стороны. Например, у вершины \(A\) соседние вершины --- \(B\) и \(F\).

\diagonalfigure

\begin{propertybox}[Диагональ]
  Отрезок, соединяющий две несоседние вершины многоугольника, называется
  \textbf{диагональю}.
\end{propertybox}

Стороны \(AB\) и \(AF\) не являются диагоналями. В шестиугольнике из
вершины \(A\) можно провести три диагонали: \(AC\), \(AD\), \(AE\).

\begin{conclusionbox}[Сколько диагоналей выходит из одной вершины]
  В \(n\)-угольнике выбранную вершину нельзя соединять диагональю с самой
  собой и с двумя соседними вершинами. Поэтому из неё выходит
  \[
    n-3
  \]
  диагонали.
\end{conclusionbox}

\section{На какие фигуры диагональ делит многоугольник}

\pentagonsplitfigure

Диагональ \(AC\) разделила пятиугольник \(ABCDE\) на две части:
треугольник \(ABC\) и четырёхугольник \(ACDE\).

Чтобы назвать получившуюся фигуру, нужно обойти её границу и перечислить
вершины подряд.

\begin{questionsbox}[Проверьте себя]
  \begin{enumerate}
    \item Чем звенья ломаной отличаются от сторон многоугольника?
    \item Как узнать, замкнута ли ломаная?
    \item Как правильно назвать многоугольник?
    \item Какие вершины называются соседними?
    \item Чем диагональ отличается от стороны?
    \item Сколько диагоналей выходит из одной вершины семиугольника?
  \end{enumerate}
\end{questionsbox}
